\chapter{Curvas en el plano y el espacio}
\section{Curvas planas y curvas regulares}
\begin{definition}[Curva]
Una \textbf{curva} (en $\displaystyle \R^{n} $) es una aplicación $\displaystyle \mathcal{C}^{\infty } $ de la forma  
\[ \alpha : I = \left(a,b\right)\to \R^{n} : t \to \alpha\left(t\right)=\left(x_{1}\left(t\right), \ldots, x_{n}\left(t\right)\right).\]
\end{definition}
\begin{notation}
	Es importante tener en cuenta:
	\begin{enumerate}
	\item Si $\displaystyle \alpha : I \to \R^{2} $ diremos que es una \textbf{curva en el plano} o una \textbf{curva plana}. Por otro lado, para $\displaystyle \alpha : I \to \R^{3} $ diremos que es una \textbf{curva en el espacio}.
	\item Diremos que $\displaystyle f $ es diferenciable si y solo si $\displaystyle f \in \mathcal{C}^{\infty } $.
	\item Utilizamos $\displaystyle I $ para designar el dominio. Lo tomamos abierto de la forma $\displaystyle \left(a,b\right) $, $\displaystyle \left(-\infty, \infty\right) $ o $\displaystyle \left(-\infty, b\right) $. No lo tomamos de la forma $\displaystyle \left[a,b\right]  $ porque pueden surgir problemas de derivabilidad.
	\item Cuando hablamos de curva hablamos de \textbf{curvas parametrizadas} y \textbf{curvas diferenciables}.
	\end{enumerate}
\end{notation}
\begin{eg}
\begin{enumerate}
\item Recta parametrizada:
	\[\alpha\left(t\right)=\left(a_{1}+a_{2}\right) + t\left(v_{1}, v_{2}\right) \Rightarrow 
	\begin{cases}
	x\left(t\right)= a_{1} + tv_{1} \\
	y\left(t\right)= a_{2} + tv_{2}
	\end{cases}
	.\]
\item Gráfica de $\displaystyle y = \sin x $, es decir, $\displaystyle \Gamma = \left\{ \left(x, \sin x\right) \; : \; x \in \R\right\}  $. La curva sería
	\[
	\begin{cases}
	x = t \\
	y = \sin t
	\end{cases}
	.\]
En general, la gráfica de la función $\displaystyle y = f\left(x\right) $ se puede parametrizar de la forma
\[
\begin{cases}
x\left(t\right)=t \\
y\left(t\right)=f\left(t\right)
\end{cases}
.\]
\end{enumerate}

\end{eg}
\begin{definition}[Traza]
	Se llama \textbf{traza de la curva} $\displaystyle \alpha : I \to \R^{2} $ al conjunto imagen $\displaystyle \alpha \left(I\right)= \left\{ \left(x\left(t\right), y\left(t\right)\right) \; : \: t \in I\right\}  $.
\end{definition}
\begin{observation}
Una misma traza puede venir de varias curvas parametrizadas. Por ejemplo, las curvas $\displaystyle \alpha\left(t\right)=\left(t, \sin t\right) $ y $\displaystyle \beta\left(t\right)=\left(t ^{3}, \sin t ^{3}\right) $ tienen la misma traza.
\end{observation}
\begin{eg}[Describir la traza de una curva]
Consideremos la curva $\displaystyle \alpha\left(t\right)=\left(t-1, t\left(t-1\right)\right) $. Para hacer un dibujo primero hacemos un dibujo de la gráfica de $\displaystyle x\left(t\right) $ y de $\displaystyle y\left(t\right) $. A partir de ahí, intentamos hacer la gráfica de $\displaystyle y $ sobre $\displaystyle x $ mediante las dos gráficas anteriores. Otra forma de hacerlo es despejando analíticamente para obtener una expresión de $\displaystyle y\left(x\right) $ o $\displaystyle x\left(y\right) $. 
\begin{center}
    \includegraphics[width=0.8\textwidth]{~/Desktop/Images/GeoDiff1.png}
\end{center}
\end{eg}
\begin{eg}
Consideremos el caso $\displaystyle \alpha\left(t\right)=\left(2,3\right) $. Tendremos que su traza es un único punto.
\end{eg}
\begin{eg}[Circunferencia]
Consideremos la circunferencia $\displaystyle x^{2}+y^{2}=1 $. Una parametrización muy común es
\[
\begin{cases}
x = \cos\theta \\
y = \sin \theta
\end{cases}
.\]
Así, consideramos la parametrización 
\[\alpha : \R \to \R^{2} : \theta \to \left(\cos \theta, \sin \theta \right) .\]
Esta parametrización no es inyectiva puesto que $\displaystyle \alpha\left(\theta\right)=\alpha\left(\theta + 2\pi k\right) $ para $\displaystyle k \in \Z $. Para obtener una parametrización inyectiva cambiamos el dominio. El problema es que tenemos que coger como dominio un intervalo abierto, pero es difícil puesto que si cogemos $\displaystyle I = \left(0, 2\pi \right) $, nos falta por considerar un punto. Sea
\[\mathbb{S}^{1}= \left\{ x^{2} +y^{2} = 1\right\}  .\]
No es posible con una parametrización inyectiva conseguir la traza de $\displaystyle \mathbb{S}^{1} $. Una forma de solucionar el problema es coger un intervalo $\displaystyle I = [0, 2\pi ) $, pero esto no es aceptado por la definición, puesto que la inversa de la parametrización no es continua. \footnote{Tendríamos que $\displaystyle \alpha  $ no es un homeomorfismo.}  
\end{eg}

\begin{definition}
Sea $\displaystyle \alpha\left(t\right) $ una curva.
\begin{enumerate}
	\item Se llama \textbf{rama infinita} cuando $\displaystyle \|\alpha\left(t\right)\| \to \infty $ si $\displaystyle t \to t_{0}^{\pm } $ con $\displaystyle t_{0} \in \R \cup \left\{ \pm \infty \right\}  $. \footnote{Todo el rato usamos la norma euclídea.} 
	\item Se llama \textbf{rama parabólica} en la dirección $\displaystyle \left(v_{1}, v_{2}\right) \in \R^{2} / \left\{ \left(0,0\right)\right\}  $ unitario si
		\[\lim_{t \to t_{0}^{\pm }}\frac{\alpha\left(t\right)}{\|\alpha\left(t\right)\|}=\left(v_{1}, v_{2}\right) .\]
	\item Se llama \textbf{asíntota} si es rama parabólica y hay otra recta $\displaystyle r \equiv \left(v_{1}t + a_{1}, v_{2}t + a_{2}\right) $ llamada \textbf{asíntota} tal que si $\displaystyle t \to t_{0}^{\pm } $,
		\[d\left(\alpha\left(t\right), r\right) \to 0 .\]
\end{enumerate}
\end{definition}
\begin{observation}
Las asíntotas son ramas parabólicas y las ramas parabólicas son ramas infinitas.
\end{observation}

\begin{eg}
La espiral $\displaystyle \alpha\left(t\right)=\left(e^{t}\cos t, e^{t}\sin t\right) $ no tiene rama parabólica porque se va aproximando al origen.
\begin{center}
    \includegraphics[width=0.7\textwidth]{~/Desktop/Images/GeoDiff2.png}
\end{center}
Por otro lado, $\displaystyle \alpha\left(t\right)= \left(t,\sin t\right) $ tiene una rama parabólica en el $\displaystyle \left(1,0\right) $.
\end{eg}
\begin{eg}
Consideremos la curva $\displaystyle xy = y^{2} + 1 $. La parametrizamos de la forma
\[
\begin{cases}
x = t + \frac{1}{t} \\
y = t
\end{cases}
,\]
con $\displaystyle t \in \left(-\infty, 0\right)\cup \left(0, \infty\right) $. Tenemos cuatro ramas infinitas:
\begin{enumerate}
\item $\displaystyle t \to -\infty  $, entonces $\displaystyle \left(x,y\right)\to \left(-\infty, -\infty \right) $. 
\item $\displaystyle t \to 0^{-} $, entonces $\displaystyle \left(x,y\right)\to \left(-\infty, 0\right) $.
\item $\displaystyle t \to 0^{+} $, entonces $\displaystyle \left(x,y\right)\to \left(\infty, 0\right) $. 
\item $\displaystyle t \to \infty $, entonces $\displaystyle \left(x,y\right)\to \left(\infty, \infty \right) $. 
\end{enumerate}
Estudiemos el caso particular de $\displaystyle t \to - \infty $ para ver el tipo de rama de infinito:
\[ \frac{\alpha \left(t\right)}{\|\alpha\left(t\right)\|}= \frac{\left(t + \frac{1}{t}, t\right)}{\|\left(t + \frac{1}{t}, t\right)\|} \to \frac{\left(-1,-1\right)}{\sqrt{2}} .\]
Así, tenemos que $\displaystyle m = 1 $ y 
\[b = \lim_{t \to -\infty}\left(y-mx\right)=\lim_{t \to -\infty}\left(t + \frac{1}{t}-t\right)=0 .\]
Así, tiene una asíntota que es $\displaystyle y - x = 0 $. Podemos comprobarlo de esta forma:
\[d\left(\alpha\left(t\right), r\right)=\frac{ \left|y\left(t\right)-x\left(t\right)\right|}{\sqrt{2}} \to 0 .\]
\begin{center}
    \includegraphics[width=0.5\textwidth]{~/Desktop/Images/GeoDiff4.png}
\end{center}
\end{eg}

\begin{definition}[Punto doble]
Se llama \textbf{punto doble (múltiple)} a $\displaystyle P = \alpha\left(t_{1}\right)=\alpha\left(t_{2}\right) $ con $\displaystyle t_{1}\neq t_{2} $. 
\end{definition}
\begin{eg}
	Veamos si $\displaystyle C = \left\{ y^{2}=x^{2}\left(x+1\right)\right\}  $ es la traza de una curva. Podemos comprobar que 
	\[y = \pm \sqrt{x^{2}\left(x+1\right)} = \pm \left|x\right|\sqrt{x+1} .\]
Intentemos parametrizar la curva. Una forma de hacerlo es para cada pendiente $\displaystyle m $ buscar la intersección
\[
\begin{cases}
y = mx \\
y^{2} = x^{2}\left(x+1\right)
\end{cases}
.\]
Obtenemos que
\[m^{2}x^{2} = x^{2}\left(x+1\right) \Rightarrow x = m^{2}-1 .\]
Así, obtenemos la parametrización
\[
\begin{cases}
x = m^{2}-1 \\
y = m\left(m^{2}-1\right)
\end{cases}
.\]
El único punto doble de la curva es $\displaystyle \alpha\left(-1\right)=\left(0,0\right)=\alpha\left(1\right) $.  
Podemos ver que 
\[\alpha\left(m\right)=\left(m^{2}-1, m^{3}-m\right), \quad \alpha'\left(m\right)=\left(2m, 3m^{2}-1\right)\neq \left(0,0\right) .\]
\begin{center}
    \includegraphics[width=0.5\textwidth]{~/Desktop/Images/GeoDiff3.png}
\end{center}
\end{eg}
\begin{eg}
	Consideremos $\displaystyle D = \left\{ y^{2}=x^{3}\right\}  $. Algebraicamente podemos escribir
	\[y = \pm x^{\frac{3}{2}} .\]
	La podemos parametrizar con $\displaystyle t \in \R $ de la forma
	\[
	\begin{cases}
	x = t ^{2} \\
	y = t ^{3}
	\end{cases}
	.\]
	Esta es una parametrización inyectiva. Al punto $\displaystyle P = \alpha\left(0\right) = \left(0,0\right) $ se le llama \textbf{pico} o \textbf{cúspide}. Podemos ver que además en este punto
	\[\alpha'\left(0\right)=\left(0,0\right) .\]
	\begin{center}
    \includegraphics[width=0.5\textwidth]{~/Desktop/Images/GeoDiff5.png}
\end{center}

\end{eg}

\begin{definition}[Vector velocidad]
Si $\displaystyle \alpha\left(t\right) $ es curva parametrizada, se llama \textbf{vector velocidad} a $\displaystyle \alpha'\left(t\right)=\frac{d\alpha }{dt} $. 
\end{definition}
\begin{definition}[Punto regular]
Decimos que $\displaystyle t_{0}\in I $ es un \textbf{punto regular} si $\displaystyle \alpha'\left(t_{0}\right)\neq 0 $. Decimos que $\displaystyle \alpha  $ es \textbf{regular} si $\displaystyle \alpha'\left(t\right)\neq 0 $, $\displaystyle \forall t \in I $. Decimos que $\displaystyle t_{0} \in I $ es un \textbf{punto singular} si $\displaystyle \alpha'\left(t_{0}\right)=0 $.
\end{definition}

\begin{eg}
	Consideremos $\displaystyle \beta\left(t\right)=\left(t ^{2}, t ^{2}\right) $. Es una parametrización no inyectiva y la traza es $\displaystyle \left\{ y = x, \; x \geq 0\right\}  $. Lo interesante de esta parametrización es ver que aunque $\displaystyle \beta'\left(0\right)=\left(0,0\right) $, no se trata de un pico.
\end{eg}
\begin{observation}
Ya vimos en el ejemplo de la curva $\displaystyle y^{2} = x^{2}\left(x + 1\right) $ que no es inyectiva y sin embargo la derivada de la parametrización que dimos no se anula nunca, por lo que es regular. 
\end{observation}
\begin{definition}[Recta tangente]
Si $\displaystyle t_{0} $ es un punto regular de $\displaystyle \alpha\left(t\right) $, se llama \textbf{recta tangente} a $\displaystyle \alpha\left(t\right) $ en $\displaystyle t = t_{0} $ a la recta $\displaystyle T = P_{0}+L\left(\alpha'\left(t_{0}\right)\right) $. 
\end{definition}
\begin{prop}
Dada un curva $\displaystyle \alpha : I \to \R^{n} $, si $\displaystyle t_{0} \in I $ es un punto regular, entonces hay una recta $\displaystyle T $ que pasa por $\displaystyle P_{0}=\alpha\left(t_{0}\right) $ tal que 
\[d\left(\alpha\left(t\right), T\right)\leq C\left(t-t_{0}\right)^{2} , \; C \in \R,\]
en un entorno de $\displaystyle t_{0} $. 
\end{prop}
\begin{proof}
	Consideremos un entorno $\displaystyle \left[t_{0}-\delta , t_{0} + \delta \right] \subset I $ con $\displaystyle \delta > 0 $ y sea $\displaystyle t \in \left[t_{0}-\delta, t_{0} + \delta \right]  $. Por el teorema de Taylor existe $\displaystyle \xi \in \left(t, t_{0}\right) $ tal que,
\[\alpha\left(t\right)=\alpha\left(t_{0}\right)+\alpha'\left(t_{0}\right)\left(t-t_{0}\right)+ \frac{\alpha''\left(\xi\right)}{2}\left(t-t_{0}\right)^{2} .\]
De esta forma, tenemos que cogiendo $\displaystyle T\left(t\right)=P_{0} + \alpha'\left(t_{0}\right)\left(t-t_{0}\right) $ como la recta tangente, 
\[
\begin{split}
	d\left(\alpha\left(t\right), T\left(t\right)\right) \leq & \left\|\left(\alpha\left(t_{0}\right)+\alpha'\left(t_{0}\right)\left(t-t_{0}\right)+\frac{\alpha''\left(\xi\right)}{2}\left(t-t_{0}\right)^{2}\right)-\left(P_{0}+\alpha'\left(t_{0}\right)\left(t-t_{0}\right)\right)\right\|\\
	= & \frac{\|\alpha''\left(\xi\right)\|}{2}\left(t-t_{0}\right)^{2} .
\end{split}
\]
Finalmente, como $\displaystyle \alpha''\left(t\right) $ es continua tendremos que está acotada en $\displaystyle \left[t_{0}-\delta, t_{0} + \delta \right]  $, por lo que existe $\displaystyle M > 0 $ tal que $\displaystyle \alpha''\left(t\right)\leq M $, $\displaystyle \forall t \in \left[t_{0}-\delta, t_{0} + \delta \right]  $ y entonces $\displaystyle d\left(\alpha\left(t\right), T\left(t\right)\right)\leq \frac{M}{2}\left(t-t_{0}\right)^{2} $.  
\end{proof}

\section{Curvas implícitas}
Sea $\displaystyle U\subset\R^{2} $ un abierto y $\displaystyle F : U \to \R $ una función $\displaystyle \mathcal{C}^{\infty} $. Consideremos el conjunto $\displaystyle C = \left\{ \left(x,y\right)\in U \; : \; F\left(x,y\right)= 0\right\}  $. 
\begin{definition}[Curva de nivel]
Las \textbf{curvas de nivel} son conjuntos de la forma $\displaystyle C_{\lambda } = \left\{ \left(x,y\right) \; : \; F\left(x,y\right)=\lambda \right\} = \left\{ \left(x,y\right) \; : \; F_{\lambda }\left(x,y\right)=F\left(x,y\right)-\lambda = 0\right\}  $. 
\end{definition}
\begin{observation}
Claramente $\displaystyle \nabla F_{\lambda }=\nabla F $.
\end{observation}

\begin{lema}
Si $\displaystyle y = h\left(x\right) $ es $\displaystyle \mathcal{C}^{\infty} $, entonces la gráfica es una curva parametrizada regular. 
\end{lema}
\begin{proof}
Tomando la parametrización $\displaystyle \alpha\left(t\right)=\left(t, h\left(t\right)\right) $, tenemos que $\displaystyle \alpha'\left(t\right)=\left(1, h'\left(t\right)\right)\neq \left(0,0\right) $, $\displaystyle \forall t \in I $.
\end{proof}
\begin{definition}[Curva implícita regular]
	Si $\displaystyle C = \left\{ \left(x,y\right)\in U \; : \; F\left(x,y\right)=0\right\}  $ y $\displaystyle \nabla F\left(x,y\right)\neq \left(0,0\right) $, $\displaystyle \forall \left(x,y\right)\in C $, llamamos a $\displaystyle C $ \textbf{curva implícita regular}.
\end{definition}
\begin{theorem}
Si $\displaystyle \left(x_{0}, y_{0}\right)\in C $ y $\displaystyle \nabla F\left(x_{0}, y_{0}\right)\neq \left(0,0\right)$, entonces existe $\displaystyle \epsilon > 0 $ tal que $\displaystyle B\left(\left(x_{0}, y_{0}\right), \epsilon \right)\cap C= C_{\epsilon} $ es la traza de una curva parametrizada regular implícita.
\end{theorem}
\begin{proof}
Tenemos que por hipótesis
\[\nabla F\left(x_{0}, y_{0}\right)=\left(\frac{\partial F}{\partial x}\left(x_{0}, y_{0}\right), \frac{\partial F}{\partial y}\left(x_{0}, y_{0}\right)\right)\neq \left(0,0\right) .\]
Supongamos sin pérdida de generalidad que $\displaystyle \frac{\partial F}{\partial y}\left(x_{0}, y_{0}\right)\neq 0 $. Consideramos la función $\displaystyle G : U \subset \R^{2}\to \R^{2} $ tal que $\displaystyle G\left(x,y\right)=\left(F\left(x,y\right), x\right) $. Como 
\[JG\left(x_{0}, y_{0}\right)=\begin{pmatrix} \frac{\partial F}{\partial x}\left(x_{0}, y_{0}\right) & \frac{\partial F}{\partial y}\left(x_{0}, y_{0}\right) \\ 1 & 0 \end{pmatrix} \Rightarrow \left|JG\left(x_{0}, y_{0}\right)\right|\neq 0 ,\]
por el teorema de la función inversa tenemos que existe $\displaystyle \epsilon > 0 $ tal que $\displaystyle G : B\left(\left(x_{0}, y_{0}\right), \epsilon \right)\to G\left(B\left(\left(x_{0}, y_{0}\right), \epsilon \right)\right) $ es un difeomorfismo. \footnote{Queremos decir que $\displaystyle G $ es diferenciable $\displaystyle \mathcal{C}^{\infty} $ y biyectiva y $\displaystyle G^{-1} $ es diferenciable $\displaystyle \mathcal{C}^{\infty} $.} Así, tenemos que 
\[\left(x,y\right)=G^{-1}\left(u,v\right)=\left(x\left(u,v\right), y\left(u,v\right)\right) .\]
Habiendo cogido $\displaystyle u = F\left(x,y\right) $ y $\displaystyle v = x $, tenemos que 
\[
\begin{split}
	C_{\epsilon }= & \left\{ \left(x,y\right) \in B_{\epsilon }\left(x_{0}, y_{0}\right)\; : \; F\left(x,y\right)=0\right\}  \\
	= & \left\{ \left(x,y\right) \; : \; y = y\left(u,v\right), u = 0, x = v\right\} \\
	= & \left\{ \left(x,y\right)\; : \; y = y\left(0,x\right)\right\} = \left\{ \left(x,y\right) \; : \; y = h\left(x\right)\right\}  ,
\end{split}
\]
tomando $\displaystyle h\left(x\right)=y\left(0,x\right) $. 
\end{proof}
\begin{eg}
	Consideremos $\displaystyle F\left(x,y\right)=x^{2}+y^{2}=1 $ y $\displaystyle C = \left\{ x^{2}+y^{2}=1\right\}  $. Tenemos que $\displaystyle \nabla F = \left(2x, 2y\right) $ y como $\displaystyle \left(0,0\right)\not\in C $ tenemos que $\displaystyle \nabla F\left(x,y\right)\neq 0 $, $\displaystyle \forall\left(x,y\right)\in C $. Tenemos que 
	% pedirle a Andrés el ejemplo completo
\end{eg}
\begin{observation}
	Recordamos que $\displaystyle \mathbb{S}^{1} $ se puede cubrir con dos parametrizaciones inyectivas, pero no con una. 
\end{observation}
\begin{observation}
Recordamos que la curva $\displaystyle y^{2}-x^{2}\left(x+1\right)=0 $ tiene gradiente nulo en el origen y sin embargo se cumplen las conclusiones del teorema.
\end{observation}

\section{Cambios de parámetro y parámetro arco}
\begin{definition}[Cambio de parámetro]
Se llama \textbf{cambio de parámetro} a una función $\displaystyle h: I \to J $; $\displaystyle I,J\subset \R $ intervalos, $\displaystyle h \in \mathcal{C}^{\infty} $ biyectiva y $\displaystyle h'\left(x\right)\neq 0 $, $\displaystyle \forall x \in I $.
\end{definition}
\begin{observation}
Como $\displaystyle h'\left(x\right)\neq 0 $, la función inversa es derivable $\displaystyle \mathcal{C}^{\infty} $ y además
\[\left(h^{-1}\right)'\left(h\left(x\right)\right)=\frac{1}{h'\left(x\right)} .\]
Es decir, los cambios de parámetros son difeomorfismos.
\end{observation}

\begin{notation}
Claramente $\displaystyle h'\left(x\right) $ tiene signo constante por ser biyectiva. Distinguimos dos casos:
\begin{itemize}
\item Si $\displaystyle h'\left(x\right) > 0 $, es estrictamente creciente (\textbf{cambio de parámetro positivo}).
\item Si $\displaystyle h'\left(x\right) < 0 $, es estrictamente decreciente (\textbf{cambio de parámetro negativo}).
\end{itemize}
\end{notation}

\begin{definition}[Reparametrización]
	Sea $\displaystyle \alpha : I \to \R^{n} $ una curva parametrizada por $\displaystyle t $. Sea $\displaystyle u = h\left(t\right) $ un cambio de parámetros. Se llama \textbf{reparametrización} de $\displaystyle \alpha \left(t\right) $ con $\displaystyle u = h\left(t\right) $ a la curva $\displaystyle \tilde{\alpha} = \alpha \circ h^{-1} : J \to \R^{n} $ con $\displaystyle \tilde{\alpha}\left(u\right)=\alpha\left(t\right)=\alpha\left(h^{-1}\left(u\right)\right) $.
\end{definition}
\begin{eg}
Consideremos la curva $\displaystyle \alpha \left(t\right)=\left(t ^{2}, t ^{3}\right) $ y consideremos el cambio de parámetro 
\[t = \sin u = h^{-1}\left(u\right) \iff u = \arcsin t = h\left(t\right) .\]
Así, tenemos que 
\[\tilde{\alpha}\left(u\right)=\alpha\left(h^{-1}\left(u\right)\right)=\alpha\left(\sin u\right)=\left(\sin ^{2}u, \sin ^{3}u\right) .\]
Normalmente se abusa de notación de esta forma: $\displaystyle \alpha\left(u\right)=\left(\sin ^{2}u, \sin ^{3}u\right) $. 
\end{eg}
\begin{observation}[Recta velocidad]
 Supongamos que $\displaystyle \tilde{\alpha} = \alpha \circ h^{-1} $, por lo que $\displaystyle \alpha\left(t\right)=\tilde{\alpha}\left(h\left(t\right)\right) $, por lo que si $\displaystyle u_{0}=h\left(t_{0}\right) $, entonces $\displaystyle \tilde{\alpha}\left(u_{0}\right)=\alpha\left(t_{0}\right) $. De esta forma, aplicando la regla de la cadena obtenemos que 
		\[\alpha'\left(t_{0}\right)=\tilde{\alpha}'\left(h\left(t_{0}\right)\right)h'\left(t_{0}\right) \Rightarrow \alpha'\left(t_{0}\right)=\lambda \tilde{\alpha}'\left(u_{0}\right) ,\]
	siendo $\displaystyle \lambda = h'\left(t_{0}\right) $. Así, son proporcionales y si $\displaystyle h'\left(t_{0}\right)>0 $ van en el mismo sentido y si $\displaystyle h'\left(t_{0}\right)<0 $ van en sentido opuesto. 
\end{observation}
\subsection*{Longitud de una curva parametrizada}
Sea $\displaystyle \alpha : I\to \R^{n} $ con $\displaystyle \left[a,b\right] \subset I $ una curva parametrizada. Calculamos la longitud entre $\displaystyle t = a $ y $\displaystyle t = b $. Dado $\displaystyle n \in \N $ tomamos la partición $\displaystyle P_{n} = \left\{ a < t_{1} < \cdots < t_{n-1}< b\right\}  $ con $\displaystyle t_{i}-t_{i-1}=\frac{b-a}{n} $ y definimos 
\[ L_{n} := \sum^{n}_{i = 1} \|\alpha\left(t_{i}\right)-\alpha\left(t_{i-1}\right)\| .\]
Una forma razonable de calcular la longitud que buscamos sería $\displaystyle \ell^{b}_{a}\left(\alpha \right)= \lim_{n \to \infty}L_{n} $. Aplicando el teorema del valor medio, $\displaystyle \forall i \in \left\{ 1, \ldots, n\right\}  $, $\displaystyle \exists\xi_{i} \in \left(t_{i-1}, t_{i}\right) $ tal que 
\[
\begin{split}
	\ell ^{b}_{a}\left(\alpha \right)= &\lim_{n \to \infty}L_{n} = \lim_{n \to \infty}\sum^{n}_{i = 1}\|\alpha\left(t_{i}\right)-\alpha\left(t_{i-1}\right)\| = \lim_{n \to \infty}\sum^{n}_{i = 1} \|\alpha'\left(\xi_{i}\right)\left(t_{i}-t_{i-1}\right)\| \\
	= &  \lim_{n \to \infty}\frac{b-a}{n}\sum^{n}_{i = 1}\|\alpha'\left(\xi_{i}\right)\| = \int^{b}_{a} \|\alpha'\left(t\right)\| \; dt  .
\end{split}
\]
\begin{definition}[Longitud de una curva]
	Sea $\displaystyle \alpha : I \to \R^{n} $ una curva y $\displaystyle \left[a,b\right] \subset I $. Se llama \textbf{longitud} de $\displaystyle \alpha \left(t\right) $ entre $\displaystyle a $ y $\displaystyle b $ a
\[\ell^{b}_{a}\left(\alpha \right):=\int^{b}_{a} \|\alpha'\left(t\right)\| \; dt .\]
\end{definition}
\begin{prop}
	Sea $\displaystyle \alpha : I \to \R^{n} $ y $\displaystyle \tilde{\alpha} = \alpha \circ h^{-1} $ su reparametrización con cambio de parámetro $\displaystyle u = h\left(t\right) $. Si $\displaystyle \left[a,b\right] \subset I $ y $\displaystyle h^{-1}\left(\left[a,b\right] \right)=\left[c,d\right]  $, entonces $\displaystyle \ell^{d}_{c}\left(\tilde{\alpha}\right)=\ell^{b}_{a}\left(\alpha \right) $. 
\end{prop}
\begin{proof}
Aplicando el teorema de cambio de variable,
\[
	\ell^{b}_{a}\left(\alpha \right)=  \int^{b}_{a} \|\alpha'\left(t\right)\| \; dt = \int^{b}_{a} \|\tilde{\alpha}'\left(h\left(t\right)\right)\| \left|h\left(t\right)\right| \; dt = \int^{d}_{c} \|\tilde{\alpha}'\left(u\right)\| \; du = \ell^{d}_{c}\left(\tilde{\alpha}\right).
\]
\end{proof}
\begin{definition}[Función espacio recorrido]
Dada una curva $\displaystyle \alpha : I \to \R^{n} $, se llama \textbf{función espacio recorrido} $\displaystyle s : I \to \R $ desde $\displaystyle t_{0} \in I $ a 
\[s\left(t\right):=\int^{t}_{t_{0}} \|\alpha'\left(u\right)\| \; d u .\]
\end{definition}
\begin{observation}
\begin{enumerate}
	\item Si $\displaystyle t > t_{0} $ y $\displaystyle \left[t_{0}, t\right] \subset I $, tendremos que $\displaystyle s\left(t\right)=\ell^{t}_{t_{0}}\left(\alpha \right)\geq 0 $. Similarmente, si $\displaystyle t < t_{0} $ y $\displaystyle \left[t, t_{0}\right] \subset I $, tendremos que $\displaystyle s\left(t\right)=-\ell^{t_{0}}_{t}\left(\alpha \right) $.
	\item Si $\displaystyle t = t_{0} $ claramente $\displaystyle s\left(t_{0}\right)=0 $.
	\item Tenemos que $\displaystyle v = \frac{d s}{dt}=\|\alpha'\left(t\right)\| $.
	\item Si tomamos $\displaystyle t_{i}\in I $, la función $\displaystyle s_{1}\left(t\right)=\int^{t}_{t_{1}} \|\alpha'\|  = \int^{t_{0}}_{t_{1}} \|\alpha' \| + \int^{t}_{t_{0}} \|\alpha'\| = c_{0} + s\left(t\right) $.
	\item También se llama a $\displaystyle s\left(t\right) $ \textbf{longitud de arco}.
\end{enumerate}
\end{observation}
\begin{lema}
Si $\displaystyle \alpha\left(t\right)  $ es una curva regular, entonces $\displaystyle s = s\left(t\right) $ es un cambio de parámetros positivo.
\end{lema}
\begin{proof}
Está claro puesto que $\displaystyle \frac{d s}{dt}= \|\alpha'\left(t\right)\| > 0 $.
\end{proof}
\begin{definition}[Reparametrización longitud de arco]
Al cambio de parámetros $\displaystyle s\left(t\right) $ se le llama parametrización \textbf{longitud de arco}. Además, $\displaystyle \alpha\left(s\right)=\alpha\left(s\left(t\right)\right) $ es la reparametrización de $\displaystyle \alpha  $ por longitud de arco (p.p.a).
\end{definition}

\begin{eg}
	Consideremos la circunferencia $\displaystyle \alpha\left(t\right)=\left(R \cos t, R \sin t\right) $ con $\displaystyle t \in \left[0, 2\pi \right]  $. Tenemos que $\displaystyle \alpha'\left(t\right)=\left(-R\sin t, R \cos t\right) $, por lo que $\displaystyle \|\alpha'\left(t\right)\| = R $ y 
	\[\ell^{2\pi }_{0}\left(\alpha \right)=\int^{2\pi }_{0} \|\alpha'\left(t\right)\| \; dt = 2\pi R .\]
Alternativamente, el parámetro arco
\[s\left(t\right)=\int^{t}_{0} \|\alpha'\| = \int^{t}_{0} R \; d u = Rt .\]
Así, $\displaystyle t = \frac{s}{R} $ y nos queda que 
\[\alpha\left(s\right)=\left(R\cos \frac{s}{R}, R\sin \frac{s}{R}\right) .\]
Esta es la circunferencia reparametrizada por longitud de arco con $\displaystyle s \in \left[0, 2\pi R\right]  $. 
\end{eg}
\begin{eg}[Cicloide]
Cogemos $\displaystyle t \in \R $ y la parametrización es
\[
\begin{cases}
x = t - \sin t \\
y = 1- \cos t
\end{cases}
.\]
Así, tenemos que la curva será $\displaystyle \alpha\left(t\right)=\left(t-t\sin t, 1-\cos t\right) $ y por tanto $\displaystyle \alpha'\left(t\right)=\left(1-\cos t, \sin t\right) $. Así, 
\[\|\alpha'\left(t\right)\| = \sqrt{\left(1-\cos t\right)^{2}+ \sin ^{2}t} = \sqrt{2 - 2 \cos t} =2\sqrt{\frac{1-\cos t}{2}} = 2 \left|\sin \frac{t}{2}\right| = 2\sin \frac{t}{2}.\]
Tenemos que $\displaystyle \|\alpha'\left(t\right)\| = 0 $ en $\displaystyle t = 2k\pi  $ con $\displaystyle k \in \Z $. Como la función es $\displaystyle 2\pi  $-periódica, nos restringimos a $\displaystyle t \in \left(0,2\pi \right) $ donde $\displaystyle \alpha\left(t\right) $ es regular. Hallamos el parámetro longitud de arco cogiendo $\displaystyle t_{0} = \pi  $:
\[s\left(t\right)=\int^{t}_{\pi} \|\alpha'\left(t\right)\| \; dt = -4\cos \frac{t}{2} .\]
Así, tenemos que para $\displaystyle t \in \left(0,2\pi \right) $, $\displaystyle s \in \left(-4,4\right) $.
Reparametrizamos por el arco:
\[\frac{s}{-4}= \cos\frac{t}{2} \iff t = 2\arccos\left(-\frac{s}{4}\right) .\]
De esta forma, 
\[\alpha\left(s\right)=\left(2\arccos\left(-\frac{s}{4}\right)+\frac{s\sqrt{16-s^{2}}}{8}, \frac{16- s^{2}}{8}\right) .\]
\end{eg}

\begin{prop}
	Sea $\displaystyle \alpha : I \to \R^{n} $ una curva. 
	\begin{enumerate}
	\item Si $\displaystyle \alpha  $ es regular, $\displaystyle \alpha\left(s\right) $ verifica que $\displaystyle \left\|\frac{d \alpha }{d s}\right\| =1$.
	\item Si $\displaystyle \alpha\left(t\right) $ verifica que $\displaystyle \left\|\frac{d\alpha }{dt}\right\| =1$, entonces $\displaystyle s = t $ es parámetro arco (salvo por una constante).
	\end{enumerate}
\end{prop}
\begin{proof}
	Supongamos que $\displaystyle \alpha : I \to \R^{n} $ es una curva.
	\begin{enumerate}
		\item Como $\displaystyle \tilde{\alpha}\left(s\right)=\alpha \circ s^{-1} $, tendremos que 
			\[ \|\tilde{\alpha}'\left(s\right)\| = \|\alpha'\left(s^{-1}\left(t\right)\right)\|\frac{1}{ \left|s'\left(s^{-1}\left(t\right)\right)\right|} = 1 .\]
\item Por otro lado, si $\displaystyle \left\|\frac{d\alpha }{dt}\right\| = 1 $, entonces
\[s\left(t\right)=\int^{t}_{t_{0}} \left\|\frac{d\alpha }{dt}\right\| = \int^{t}_{t_{0}} 1 =t-t_{0} .\]
	\end{enumerate}
\end{proof}
\begin{observation}
Con el parámetro arco, para cualquier curva se tiene que $\displaystyle \ell^{b}_{a}\left(\alpha\left(s\right)\right) = b-a $. 
\end{observation}
\section{Diedro de Frenet y curvatura}
\begin{notation}
	La matriz del giro de $\displaystyle \frac{\pi }{2} $ en sentido positivo es $\displaystyle J = \begin{pmatrix} 0 & -1 \\ 1 & 0 \end{pmatrix} $. De esta forma, para $\displaystyle \left(a,b\right)\in \R^{2} $, 
	\[J\begin{pmatrix} a \\ b \end{pmatrix}=\begin{pmatrix} 0 & -1 \\ 1 & 0 \end{pmatrix}\begin{pmatrix} a \\b \end{pmatrix} = \begin{pmatrix} b \\ - a \end{pmatrix} .\]
	Así, tenemos que $\displaystyle \left\{ \begin{pmatrix} a \\ b \end{pmatrix}, \begin{pmatrix} -b \\ a \end{pmatrix}\right\}  $ es una base positiva \footnote{Esto quiere decir que el determinante es positivo.}.
\end{notation}

\begin{definition}[Diedro de Frenet]
Sea $\displaystyle \alpha : I \to \R^{2} $ una curva regular y $\displaystyle t_{0} \in I $.
\begin{itemize}
\item Se llama \textbf{vector tangente} a $\displaystyle \alpha  $ en $\displaystyle t_{0} $ a
	\[\hat{t}_{\alpha }\left(t_{0}\right) = \frac{\alpha'\left(t_{0}\right)}{\|\alpha'\left(t_{0}\right)\|} .\]
\item Se llama \textbf{vector normal} a $\displaystyle \alpha  $ en $\displaystyle t_{0} $ a
	\[\hat{n}_{\alpha }\left(t_{0}\right)=J \hat{t}_{\alpha }\left(t_{0}\right) .\]
\end{itemize}
Se llama \textbf{diedro de Frénet} a $\displaystyle \left\{ \hat{t}_{\alpha }\left(t\right), \hat{n}_{\alpha }\left(t\right)\right\}  $. 
\end{definition}

\begin{notation}
Diedro significa base ortonormal positiva. 
\end{notation}
\begin{eg}[Cicloide]
	Tendremos que el vector tangente a $\displaystyle \alpha  $ en $\displaystyle t_{0} \in I $ será
	\[\hat{t}_{\alpha }\left(t_{0}\right)=\frac{\left(1-\cos t_{0}, \sin t_{0}\right)}{2\sin \frac{t_{0}}{2}} = \left(\sin \frac{t_{0}}{2}, \cos \frac{t_{0}}{2}\right) .\]
\end{eg}
\begin{eg}
Consideremos 
\[\alpha\left(s\right)=\left(2\arccos\left(-\frac{s}{4}\right) + \frac{s\sqrt{16-s^{2}}}{8}, \frac{16-s^{2}}{8}\right) .\]
Tendremos que 
\[\hat{t}_{\alpha }= \alpha'\left(s\right)= \text{calcular} .\]
Alternativametne, podemos calcular $\displaystyle \alpha'\left(t\right) $ y $\displaystyle \|\alpha'\left(t\right)\| $ y calcular de esta forma $\displaystyle \hat{t}_{\alpha }\left(t\right) $. 
\end{eg}
\begin{lema}
Dada una curva regular $\displaystyle \alpha : I \to \R^{2} $, los vectores $\displaystyle \hat{t}_{\alpha } $ y $\displaystyle \hat{n}_{\alpha } $ no dependen de reparametrizaciones positivas.
\end{lema}
\begin{proof}
	Sea $\displaystyle u = h\left(t\right) $ un cambio de parámetros positivo y sea $\displaystyle \tilde{\alpha} = \alpha \circ h^{-1} $ la reparametrización correspondiente. Tendremos que 
	\[\hat{t}_{\alpha }\left(t\right)=\frac{\alpha'\left(t\right)}{\|\alpha'\left(t\right)\|}=\frac{\tilde{\alpha}'\left(h\left(t\right)\right) h'\left(t\right)}{\|\tilde{\alpha}'\left(h\left(t\right)\right) \|h'\left(t\right)} = \hat{t}_{\tilde{\alpha}}\left(u\right) .\]
	De aquí se sigue la independencia de $\displaystyle \hat{n}_{\alpha } $ respecto de reparametrizaciones puesto que $\displaystyle \hat{n}_{\alpha } = J\hat{t}_{\alpha } $. 
\end{proof}

\begin{prop}
Si $\displaystyle \alpha : I \to \R^{2} $ es una curva regular, $\displaystyle \hat{t}_{\alpha }\left(t\right)= \frac{d\alpha }{d s} $. 
\end{prop}
\begin{proof}
	Sabemos que si $\displaystyle \alpha : I \to \R^{2} $ es una curva regular, entonces se puede parametrizar por longitud de arco. Así, podemos coger $\displaystyle \tilde{\alpha}\left(s\left(t\right)\right)=\alpha\left(t\right) $, por lo que 
	\[\hat{t}_{\alpha }\left(t\right)=\frac{\alpha'\left(t\right)}{\|\alpha'\left(t\right)\|}=\frac{\tilde{\alpha}'\left(s\left(t\right)\right)s'\left(t\right)}{\|\tilde{\alpha}'\left(s\left(t\right)\right)\| s'\left(t\right)} = \frac{\tilde{\alpha}'\left(s\left(t\right)\right)}{\|\tilde{\alpha}'\left(s\left(t\right)\right)\|} .\]
	Como $\displaystyle \|\tilde{\alpha}'\left(s\left(t\right)\right)\| = 1 $ tendremos que $\displaystyle \hat{t}_{\alpha }\left(t\right)=\tilde{\alpha}'\left(s\left(t\right)\right) $. Otra forma de ver la demostración es con la notación de diferenciales:
	\[\hat{t}_{\alpha }\left(t\right)=\frac{\frac{d\alpha }{dt}}{\left\|\frac{d\alpha }{dt}\right\|}= \frac{\frac{d\alpha }{d s} \cdot \frac{d s}{dt}}{\left\|\frac{d\alpha }{d s}\right\|\frac{d s}{dt}} = \frac{\frac{d\alpha }{d s}}{\left\|\frac{d\alpha }{d s}\right\|}= \frac{d\alpha }{d s}.\]
\end{proof}

\begin{observation}
Si $\displaystyle \alpha \left(s\right) $ está p.p.a entonces $\displaystyle \hat{t}_{\alpha }\left(s\right)=\alpha'\left(s\right) $. 
\end{observation}

\begin{lema}
Sea $\displaystyle \hat{v} : J \subset \R \to \R^{n} : \lambda \to \hat{v}\left(\lambda\right)$. Existe $\displaystyle C \geq 0 $ tal que $\displaystyle \|\hat{v}\left(\lambda \right)\| \equiv C $ si y sólo si $\displaystyle \vec{v}'\left(\lambda \right)\perp \vec{v}\left(\lambda \right) $ \footnote{Estamos asumiendo que $\displaystyle \hat{v} $ ni su derivada son idénticamente nulas.} . 
\end{lema}
\begin{proof}
Tenemos que $\displaystyle \|\vec{v}\left(\lambda \right)\|^{2} = C $, es decir, $\displaystyle \left\langle \vec{v}\left(\lambda \right), \vec{v}\left(\lambda \right) \right\rangle = C $. Por tanto, 
\[\frac{d}{d\lambda }\left\langle \vec{v}\left(\lambda \right), \vec{v}\left(\lambda \right) \right\rangle = \left\langle \frac{d\vec{v}}{d\lambda }, \vec{v} \right\rangle +\left\langle \vec{v},\frac{d\vec{v}}{d\lambda}  \right\rangle =  0 \iff \frac{d\vec{v}}{d\lambda }\perp \vec{v} .\]
\end{proof}
\begin{colorary}
Dada una curva regular $\displaystyle \alpha : I \to \R^{2} $, se tiene que $\displaystyle \frac{d\hat{t}_{\alpha }}{dt}\perp \hat{t}_{\alpha } $.
\end{colorary}
\begin{proof}
Trivial a partir del lema anterior puesto que $\displaystyle \|\hat{t}_{\alpha }\| \equiv 1 $. 
\end{proof}
\begin{definition}[Curvatura]
Dada una curva regular $\displaystyle \alpha : I \to \R^{2} $ p.p.a, se llama \textbf{curvatura} de $\displaystyle \alpha  $ en $\displaystyle s_{0} $ a $\displaystyle \kappa_{\alpha }\left(s_{0}\right) $ tal que
\[\alpha''\left(s_{0}\right)=\frac{d\hat{t}_{\alpha }}{d s}|_{s_{0}} = \kappa_{\alpha }\left(s_{0}\right) \hat{n}_{\alpha }\left(s_{0}\right) .\]
\end{definition}

\begin{prop}
Sea $\displaystyle \alpha : I \to \R^{2} $ una curva regular.
\begin{enumerate}
\item $\displaystyle \kappa_{\alpha }=\left\langle \frac{d\hat{t}_{\alpha }}{d s}, \hat{n}_{\alpha } \right\rangle $. 
\item $\displaystyle \kappa_{\alpha }\hat{n}_{\alpha }=\alpha''\left(s\right)$ y $ \|\kappa_{\alpha }\hat{n}_{\alpha }\|= \left|\kappa_{\alpha}\right| = \|\alpha''\left(s\right)\| $.
\item $\displaystyle \kappa_{\alpha } \equiv 0 $ si y sólo si $\displaystyle \alpha\left(s\right) $ es una recta.
\end{enumerate}
\end{prop}
\begin{proof}
Sea $\displaystyle \alpha : I \to \R^{2} $ una curva regular.
\begin{enumerate}
\item Como $\displaystyle \frac{d\hat{t}_{\alpha }}{d s}= \kappa_{\alpha }\hat{n}_{\alpha } $, tenemos que 
	\[\left\langle \frac{d\hat{t}_{\alpha }}{d s}, \hat{n}_{\alpha } \right\rangle =\left\langle \kappa_{\alpha }\hat{n}_{\alpha }, \hat{n}_{\alpha } \right\rangle =\kappa_{\alpha } .\]
\item Si $\displaystyle \alpha\left(s\right) $ está p.p.a, entonces $\displaystyle \hat{t}_{\alpha }\left(s\right)=\alpha'\left(s\right) $ y $\displaystyle \hat{t}_{\alpha}'\left(s\right)=\frac{d\hat{t}_{\alpha }}{d s}=\alpha''\left(s\right) $.
\item Tenemos que $\displaystyle \kappa_{\alpha }\equiv 0 $ si y solo si $\displaystyle \alpha''\left(s\right)\equiv 0 $ si y solo si $\displaystyle \alpha'\left(s\right)=v_{0} $ si y solo si $\displaystyle \alpha\left(s\right)=P_{0} + v_{0}s $, con $\displaystyle \|v_{0}\| = 1 $.
\end{enumerate}
\end{proof}
\begin{observation}
En la base canónica, $\displaystyle \hat{t}_{\alpha }\left(s\right)=\left(\cos\theta\left(s\right), \sin\theta\left(s\right)\right) $, siendo $\displaystyle \theta\left(s\right) $ el ángulo que forman los vectores $\displaystyle \hat{t}_{\alpha }\left(s\right) $ y $\displaystyle \vec{e}_{1} $. Así, tendremos que 
	\[\frac{d\hat{t}_{\alpha }}{d s}=\left(-\sin\theta\left(s\right)\theta'\left(s\right), \cos\theta\left(s\right)\theta'\left(s\right)\right) = \theta'\left(s\right)\hat{n}_{\alpha }\left(s\right) .\]
Por tanto, 
\[\kappa_{\alpha }\left(s\right)=\frac{d\theta }{d s} .\]
Así, si $\displaystyle \frac{d\theta }{d s} > 0 $, tenemos que $\displaystyle \theta\left(s\right) $ es creciente y $\displaystyle \kappa_{\alpha }\left(s\right) > 0 $. Geométricamente, supondría que el cuerpo gira a la izquierda. Análogamente, si $\displaystyle \frac{ d\theta }{d s} < 0 $ tenemos que $\displaystyle \theta\left(s\right) $ es decreciente y $\displaystyle \kappa_{\alpha }\left(s\right) < 0 $. En este caso, el cuerpo gira a la derecha. 
\end{observation}

\begin{observation}[Parametrización con signo opuesto]
Si reparametrizamos en sentido opuesto cogiendo $\displaystyle \beta\left(s\right)=\alpha\left(-s\right) $, tenemos que 
\[
\begin{cases}
\hat{t}_{\beta }\left(s\right) = -\hat{t}_{\alpha }\left(-s\right) \\
\hat{t}_{\beta }\left(s\right)=J\hat{t}_{\beta }=-J\hat{t}_{\alpha }=-\hat{n}_{\alpha }\left(s\right)
\end{cases}
.\]
Así, tenemos que 
\[\beta''\left(s\right)=\alpha''\left(-s\right)=\kappa_{\alpha }\left(s\right)\hat{n}_{\alpha }\left(-s\right)=-\kappa_{\alpha }\left(-s\right)\hat{n}_{\beta }\left(s\right) .\]
De esta forma, $\displaystyle \kappa_{\beta }\left(s\right)=-\kappa_{\alpha }\left(-s\right) $. Así, si cambiamos el signo obtenemos la curvatura opuesta.
\end{observation}

\begin{observation}[Curvatura respecto de $\displaystyle t $]
Es sencillo ver que 
\[\frac{d\hat{t}_{\alpha }}{dt} = \frac{d\hat{t}_{\alpha }}{d s}\frac{d s}{dt}=v\kappa_{\alpha }\hat{n}_{\alpha }, \; v = \frac{d s}{dt}= \left\|\frac{d\alpha }{dt}\right\| .\]
\end{observation}

\begin{eg}[Cicloide]
Sea $\displaystyle \alpha\left(t\right)=\left(t-\sin t, 1-\cos t\right) $. Tenemos que 
\[\hat{t}_{\alpha }=\left(\sin \frac{t}{2}, \cos \frac{t}{2}\right), \quad \alpha'\left(t\right)=\left(1-\cos t, \sin t\right), \; v = \|\alpha'\left(t\right)\| = 2 \sin t .\]
Tendremos que $\displaystyle \hat{n}_{\alpha }=\left(\cos \frac{t}{2}, -\sin \frac{t}{2}\right) $. Por otro lado, 
\[\frac{d\hat{t}_{\alpha }}{dt}=v\kappa_{\alpha}\hat{n}_{\alpha} .\]
\[\left(\frac{1}{2}\cos\frac{t}{2}, -\frac{1}{2}\sin \frac{t}{2}\right)=2\sin \frac{t}{2}\kappa_{\alpha }\left(\cos \frac{t}{2}, \sin \frac{t}{2}\right) \Rightarrow \frac{1}{2} = 2\sin \frac{t}{2}\kappa_{\alpha } \Rightarrow \kappa_{\alpha }\left(t\right)=-\frac{1}{4\sin \frac{t}{2}} .\]
Si lo hacemos con la longitud de arco, 
\[\cos \frac{t}{2} = - \frac{s}{4}, \quad \sin \frac{t}{2} = \sqrt{1-\left(\frac{s}{4}\right)^{2}} = \frac{\sqrt{16-s^{2}}}{4}, \quad \kappa_{\alpha }=\frac{1}{\sqrt{16-s^{2}}} .\]
\end{eg}
\begin{theorem}[Teorema de Frénet]
	Sea $\displaystyle \alpha : I \to \R^{n} $ una curva regular. Entonces, 
\[\frac{d\hat{t}_{\alpha }}{d s}= \kappa_{\alpha }\hat{n}_{\alpha} \quad \text{y} \quad
\frac{d\hat{n}_{\alpha }}{d s}= -\kappa_{\alpha \hat{t}_{\alpha}} .\]
Equivalentemente, 
\[\frac{d\hat{t}_{\alpha }}{dt}= v\kappa_{\alpha }\hat{n}_{\alpha } \quad \text{y} \quad \frac{d\hat{n}_{\alpha }}{dt}= - v\kappa_{\alpha }\hat{t}_{\alpha } .\]
\end{theorem}
\begin{proof}
La primera igualdad se sigue por definición. Por otro lado, 
\[\frac{d\hat{n}_{\alpha }}{d s}= \frac{d\left(J\hat{t}_{\alpha }\right)}{d s}=J\left(\frac{d\hat{t}_{\alpha }}{d s}\right) =J\left(\kappa_{\alpha }\right)=\kappa_{\alpha }\left(J\hat{n}_{\alpha }\right) =\kappa_{\alpha }\left(-\hat{t}_{\alpha }\right)=-\kappa_{\alpha }\hat{t}_{\alpha}.\]
\end{proof}

\begin{observation}
Se puede usar la regla pnemotécnica:
\[\begin{pmatrix} \hat{t}_{\alpha }'\left(s\right) \\ \hat{n}_{\alpha }'\left(s\right) \end{pmatrix}=\begin{pmatrix} 0 & \kappa_{\alpha } \\-\kappa_{\alpha } & 0 \end{pmatrix}\begin{pmatrix} \hat{t}_{\alpha }\\\hat{n}_{\alpha } \end{pmatrix} .\]
\end{observation}

\begin{prop}
	Sea $\displaystyle \alpha : I \to \R^{2} $ una curva regular. Entonces,
	$\displaystyle \alpha\left(t\right) $ es un trozo de circunferencia si y solo si $\displaystyle \kappa_{\alpha }\left(t\right)\equiv k_{0} $ para algún $\displaystyle k_{0} \in \R/ \left\{ 0\right\}  $. 
\end{prop}
\begin{proof}
Vamos a reparametrizar por arco. Tenemos que $\displaystyle \alpha\left(s\right) $ es circunferencia si y solo si $\displaystyle \kappa_{\alpha }\equiv k_{0} $, es decir, $\displaystyle \frac{d\hat{t}_{\alpha }}{d s}\equiv k_{0} $ es lo que queremos demostrar.
\begin{description}
\item[($\displaystyle \Rightarrow $)] Supongamos que $\displaystyle \alpha \left(s\right) $ es una circunferencia centrada en el origen de radio $\displaystyle R $. Tenemos que 
	\[\alpha\left(s\right)=\left(R\cos \frac{s}{R}, R \sin \frac{s}{R}\right) \Rightarrow
	\begin{cases}
	\hat{t}_{\alpha }\left(s\right)=\alpha'\left(s\right)=\left(-\sin \frac{s}{R}, \cos \frac{s}{R}\right) \\
	\hat{n}_{\alpha }=\left(-\cos\frac{s}{R}, -\sin \frac{s}{R}\right)
	\end{cases}
	.\]
	Así, nos queda que 
	\[\alpha''\left(s\right)=\left(-\frac{1}{R}\cos \frac{s}{R}, -\frac{1}{R}\sin \frac{s}{R}\right)\Rightarrow \kappa_{\alpha }\left(s\right)= \frac{1}{R} .\]
Ya vimos que $\displaystyle \kappa_{\alpha }\left(-s\right)=-\frac{1}{R} $. En el primer caso, el centro de la circunferencia viene dado por
\[C_{0}= \alpha\left(s\right) + R \hat{n}\left(s\right)= \alpha\left(s\right)+\frac{1}{\kappa_{\alpha }}\hat{n}\left(s\right) .\]
En el segundo caso, 
\[C_{0}=\alpha\left(s\right)-R\hat{n}\left(s\right)=\alpha\left(s\right)+\frac{1}{\kappa_{\alpha }}\hat{n}_{\alpha }\left(s\right) .\]
\item[($\displaystyle \Leftarrow $)] Defimimos $\displaystyle C\left(s\right)=\alpha\left(s\right)+\frac{1}{\kappa_{\alpha}}\hat{n}_{\alpha }\left(s\right) $ y veamos que $\displaystyle C\left(s\right)=C_{0} $:
	\[C'\left(s\right)=\alpha'\left(s\right)+\frac{1}{\kappa_{\alpha}}\hat{n}_{\alpha }'=\hat{t}_{\alpha }+\frac{1}{\kappa_{\alpha }}\left(-\kappa_{\alpha }\hat{t}_{\alpha }\right)=0 .\]
	Así, tenemos que $\displaystyle C\left(s\right) $ es fijo. Ahora, queremos ver que $\displaystyle \|\alpha\left(s\right)-C_{0}\| $ es constante. Tenemos que
	\[ \alpha\left(s\right)-C_{0}=-\frac{1}{\kappa_{\alpha }}\hat{n}_{\alpha } .\]
De esta forma, 
\[\|\alpha\left(s\right)-C_{0}\|=\sqrt{\frac{1}{\kappa_{\alpha }^{2}}} = \frac{1}{ \left|\kappa_{\alpha }\right|} .\]
Tenemos así que $\displaystyle R = \frac{1}{ \left|\kappa_{\alpha }\right|} $. 
\end{description}
\end{proof}
\begin{definition}[Centro de curvatura y circunferencia osculadora]
Dada una curva regular $\displaystyle \alpha: I \to \R^{2} $, decimos que el \textbf{centro de curvatura} de $\displaystyle \alpha\left(s\right) $ es el punto
\[C\left(s\right):=\alpha\left(s\right)+\frac{1}{\kappa_{\alpha }\left(s\right)}\hat{n}_{\alpha }\left(s\right) .\]
La \textbf{circunferencia osculadora} es la circunferencia de centro $\displaystyle C\left(s_{0}\right) $ y radio $\displaystyle R = \frac{1}{ \left|\kappa_{\alpha }\left(s_{0}\right)\right|} $.
\end{definition}
\begin{prop}
Si $\displaystyle \alpha\left(s\right) $ es regular con $\displaystyle \kappa_{\alpha }\left(s_{0}\right)\neq 0 $ y $\displaystyle \beta\left(s\right) $ es la circunferencia con centro $\displaystyle C_{0}=C\left(s_{0}\right) $ y radio $\displaystyle R = \frac{1}{ \left|\kappa_{\alpha }\left(s_{0}\right)\right|} $ (ajustamos $\displaystyle \beta\left(s_{0}\right)=\alpha\left(s_{0}\right)=P_{0} $, con el mismo sentido de giro en $\displaystyle \alpha  $ y $\displaystyle \beta  $).
Entonces, existe $\displaystyle C \geq 0 $ tal que $\displaystyle \|\alpha\left(s\right)-\beta\left(s\right)\|\leq C \left|s-s_{0}\right|^{3} $. 
\end{prop}
\begin{proof}
Sea $\displaystyle \alpha\left(s_{0}\right)=\beta\left(s_{0}\right)=P_{0} $. Sea $\displaystyle \beta  $ la circunferencia desde $\displaystyle C_{0} $ a $\displaystyle P_{0} $, por lo que $\displaystyle \hat{n}_{\beta } || \overrightarrow{C_{0}P_{0}} $. Por otro lado, 
\[\overrightarrow{C_{0}P_{0}}=\alpha\left(s_{0}\right)-C_{0}=\frac{1}{ \left|\kappa_{\alpha }\right|}\hat{n}_{\alpha } ||\hat{n}_{\alpha } .\]
Así, $\displaystyle \hat{n}_{\beta }\left(s_{0}\right) || \hat{n}_{\alpha }\left(s_{0}\right) $. Por tanto, $\displaystyle \hat{t}_{\alpha }\left(s_{0}\right) || \hat{t}_{\beta }\left(s_{0}\right) $. Como tenemos el mismo sentido de giro debe ser que $\displaystyle \alpha'\left(s_{0}\right)=\beta'\left(s_{0}\right) $, 
\[\alpha''\left(s_{0}\right)=\kappa_{\alpha }\left(s_{0}\right)\hat{n}_{\alpha }\left(s_{0}\right) = \pm \frac{1}{R}\hat{n}_{\alpha }\left(s_{0}\right) = \kappa_{\beta }\left(s_{0}\right)\hat{n}_{\beta }\left(s_{0}\right)=\beta''\left(s_{0}\right) .\]
Aplicando el teorema de Taylor, 
\[\alpha\left(s\right) = \alpha\left(s_{0}\right) + \alpha'\left(s_{0}\right)\left(s-s_{0}\right) + \frac{\alpha''\left(s_{0}\right)}{2} \left(s-s_{0}\right)^{2} + \frac{\alpha'''\left(\xi\right)}{3!}\left(s-s_{0}\right)^{3} .\]
\[\beta\left(s\right) = \beta\left(s_{0}\right)+\beta'\left(s_{0}\right)\left(s-s_{0}\right) + \frac{\beta''\left(s_{0}\right)}{2!}\left(s-s_{0}\right)^{2} + \frac{\beta'''\left(\xi\right)}{3!} \left(s-s_{0}\right)^{3} .\]

En un intervalo $\displaystyle \left[s_{0}-\epsilon, s_{0} + \epsilon \right]  $ tenemos que $\displaystyle \|\alpha'''\| $ está acotado y $\displaystyle \|\beta'''\| $ también lo esta. Por tanto, 
\[ \|\alpha\left(s\right)-\beta\left(s\right)\| \leq \frac{C\left(\|\alpha'''\| + \|\beta'''\|\right)}{6} \left|s-s_{0}\right|^{3} .\]
\end{proof}
\begin{theorem}
Dada una curva regular $\displaystyle \alpha : I \to \R^{2} $, tenemos que 
\[\kappa_{\alpha }\left(t\right)=\frac{\det\left(\alpha', \alpha''\right)}{\|\alpha'\|^{3}} .\]
\end{theorem}
\begin{proof}
Es sencillo ver que 
\[\det\left(\frac{d\alpha }{ dt}, \frac{d^{2}\alpha }{dt ^{2}}\right) = \det\left(v\hat{t}_{\alpha }, \frac{dv}{dt}\hat{t}_{\alpha } + \kappa_{\alpha }v^{2}\hat{n}_{\alpha }\right) = \begin{vmatrix} v & \frac{dv}{dt} \\ 0 & \kappa_{\alpha }v^{2} \end{vmatrix} =\kappa_{\alpha }v^{2} .\]
Tenemos que $\displaystyle \kappa_{\alpha }=\frac{1}{v^{3}}\det\left(\frac{d\alpha }{dt}, \frac{d^{2}\alpha }{dt ^{2}}\right) $.
\end{proof}

\begin{observation}
En el caso $\displaystyle t = s $ tenemos que como $\displaystyle \|\alpha'\left(s\right)\| = 1 $,
\[\frac{\det\left(\frac{d\alpha }{d s}, \frac{d ^{2}\alpha }{d s^{2}}\right)}{\|\alpha'\left(s\right)\|^{3}} = \det\left(\hat{t}_{\alpha }, \kappa_{\alpha }\hat{n}_{\alpha }\right)=\kappa_{\alpha } .\]
\end{observation}

\begin{eg}
Dada $\displaystyle \alpha\left(t\right)= \left(t, t^{2}\right)$, tenemos que $\displaystyle \alpha'\left(t\right)=\left(1, 2t\right) $ y $\displaystyle \alpha''\left(t\right)=\left(0,2\right) $. De esta forma, 
\[\|\alpha'\left(t\right)\| = \sqrt{1 + 4t ^{2}} .\]
Por tanto, 
\[\kappa_{\alpha } = \frac{\det\left(\alpha', \alpha''\right)}{\|\alpha'\|^{3}} = \frac{1}{\left(\sqrt{1 + 4 t ^{2}}\right)}\begin{vmatrix} 1 & 0 \\ 2t & 2 \end{vmatrix} =\frac{2}{\left(1 + 4t ^{2}\right)^{\frac{3}{2}}} .\]
Hacer esto con la longitud de arco es horrible, esta es la utilidad del ejemplo. En efecto, 
\[s = \int \|\alpha'\left(t\right)\| \; dt = \int \sqrt{1 + 4t ^{2}} \; dt = \frac{2t\sqrt{1+4t ^{2}}}{4} + \frac{\log\left(2t + \sqrt{1 + 4t ^{2}}\right)}{4} .\]
Podemos ver que no se puede despejar la $\displaystyle t $. Podemos calcular $\displaystyle \alpha'\left(s\right) $ de la siguiente forma
\[\alpha'\left(s\right)=\frac{d\alpha }{d s} = \frac{d\alpha }{d t}\frac{dt}{d s} = \left(1, 2t\right)\frac{1}{\sqrt{1 + 4t ^{2}}} .\]
Análogamente, 
\[\alpha''\left(s\right)=\frac{d}{d s}\left(\frac{d\alpha }{d s}\right) = \frac{d}{dt}\left(\frac{d\alpha }{d s}\frac{dt}{d s}\right)=\frac{\left(-4t, 2\right)}{\left(1+4t ^{2}\right)^{2}} = \kappa_{\alpha }\hat{n}_{\alpha } .\]
Así, nos queda que 
\[\kappa_{\alpha }\hat{n}_{\alpha }=\kappa_{\alpha }\frac{\left(-2t, 1\right)}{\left(1 + 4t ^{2}\right)^{\frac{1}{2}}} \Rightarrow \kappa_{\alpha }=\frac{2}{\left(1 + 4t ^{2}\right)^{\frac{3}{2}}} .\]
\end{eg}
\begin{eg}[Cicloide]
	Calculemos la curvatura de la cicloide. Recordamos que $\displaystyle \alpha\left(t\right)=\left(t-\sin t, 1- \cos t\right) $ y $\displaystyle \|\alpha'\| = 2\sin \frac{t}{2} $. Además, $\displaystyle s = -4\cos \frac{t}{2} $ por lo que 
	\[\alpha\left(s\right)=\left(2\arccos\left(-\frac{s}{4}\right)+ s\frac{\sqrt{16-s ^{2}}}{8}, \frac{16-s^{2}}{8}\right) .\]
Ya habíamos visto que $\displaystyle \hat{t}_{\alpha }=\left(\sin \frac{t}{2}, \cos \frac{t}{2}\right) $ y $\displaystyle \alpha'\left(s\right)=\left(\frac{\sqrt{16-s^{2}}}{4}, - \frac{s}{4}\right) $. Por un lado, 
\[\hat{t}_{\alpha }\left(s\right)=\alpha'\left(s\right)=\left(\frac{16-s^{2}}{4}, - \frac{s}{4}\right) .\]
\[\hat{n}_{\alpha }\left(s\right)=\left(-\frac{s}{4}, -\frac{\sqrt{16-s^{2}}}{4}\right) .\]
\[\alpha''\left(s\right)=\left(-\frac{2s}{24\sqrt{16-s^{2}}}, -\frac{1}{4}\right)=\kappa_{\alpha }\hat{n}_{\alpha } .\]
Por otro lado, 
\[\kappa_{\alpha }=\frac{\det\left(\alpha'\left(t\right), \alpha''\left(t\right)\right)}{\|\alpha'\left(t\right)\|^{3}}=-\frac{1}{4\sin \frac{t}{2}} .\]
\end{eg}
\section{Curvas congruentes y teorema de congruencia}
\begin{definition}[Isometría lineal]
Se define \textbf{isometría lineal} como una aplicación $\displaystyle A : \R^{n} \to \R^{n} $ con $\displaystyle A \in \mathcal{M}_{n \times n}\left(\R\right) $ tal que $\displaystyle A^{t}A = I $ y $\displaystyle A^{-1}=A^{t} $. Además verifican 
\[\left\langle AX, AY \right\rangle =\left(AX\right)^{t}AY =X^{t}A^{t}AY = X^{t}Y = \left\langle X, Y \right\rangle .\]
Es decir, preserva el producto escalar. Además $\displaystyle A $ es una matriz ortonormal y $\displaystyle \det\left(A\right)=1 $. Si $\displaystyle \det\left(A\right)=-1 $ decimos que es una isometría negativa. 
\end{definition}
\begin{eg}
En $\displaystyle \R^{2} $, una isometría positiva puede ser un giro de ángulo $\displaystyle \theta $ y tiene la forma
\[A = \begin{pmatrix}  \cos \theta & -\sin \theta \\ \sin \theta & \cos \theta \end{pmatrix} .\]
Análogamente, el giro en sentido contrario
\[A =\begin{pmatrix} \cos \theta & \sin \theta \\ \sin \theta & -\cos\theta  \end{pmatrix} .\]
\end{eg}
\begin{definition}[Movimiento rígido]
Son las isometrías afines. Es decir, son aplicaciones de la forma $\displaystyle \varphi : \R^{n} \to \R^{n} : X \to AX + B $ con $\displaystyle A \in O_{n} $ y $\displaystyle B \in \R^{n} $. En forma matricial son de la forma
\[\begin{pmatrix} 1 & 0 \\ B & A \end{pmatrix} .\]
Las distancias se preservan:
\[ .\]
\[
\begin{split}
	d\left(\varphi\left(X\right), \varphi\left(A\right)\right) = & \|\varphi\left(X\right)-\varphi\left(Y\right)\| \\
	= &  \|\left(AX + B\right)-\left(AY + B\right)\| = \|A\left(X-Y\right)\| = \|X-Y\| = d\left(X,Y\right).
\end{split}
\]
Hay movimiento rígidos positivos y negativos. 
\end{definition}
\begin{definition}[Curvas congruentes]
Dos curvas $\displaystyle \alpha, \beta : J \to \R^{2} $ son \textbf{congruentes} si existe $\displaystyle \varphi:\R^{2} \to \R^{2} $ movimiento rígido (pos/neg) con $\displaystyle \beta\left(t\right)=\varphi\left(\alpha\left(t\right)\right) $. 
\end{definition}

